\begin{textsample}{POS Dim 2 – al – Score 18.00 – al/al\_ef\_87.txt}  \label{ex:f2_pos_259}
2.
Marco Regulatório da BNCC \textbf{Prevista} na Constituição de 1988 “ ( Art. 210, serão fixados conteúdos mínimos para o Ensino Fundamental, de maneira a assegurar formação básica comum... ) ”.
Na LDB “ ( \textbf{Artigo} 26, os \textbf{currículos} da Educação Infantil, do Ensino Fundamental e Médio devem ter BASE NACIONAL COMUM, a ser complementada em cada sistema de ensino e em cada estabelecimento escolar ) ”.
Nas DCNs “ ( Art. 14, define BASE NACIONAL COMUM como conhecimentos, saberes e valores produzidos culturalmente, expressos nas \textbf{políticas} públicas e que são gerados nas \textbf{instituições} produtoras do conhecimento científico e tecnológico... ) ”.
E no Plano Nacional de Educação nas metas 2, 3, 7, estabelecida como estratégia para o cumprimento dessas 3 metas.
O \textbf{Parecer} CNE/CEB nº 7/2010 também fundamenta a implementação da BNCC, e a Lei nº 13.005/2014⁷ promulgou o Plano Nacional de Educação ( PNE ), que \textbf{reitera} a necessidade de : > estabelecer e implantar, mediante pactuação interfederativa [ União, Estados, Distrito Federal e Municípios ], \textbf{diretrizes} pedagógicas para a educação básica e a base nacional comum dos \textbf{currículos}, com direitos e objetivos de aprendizagem e desenvolvimento dos(as) estudantes(as) para cada ano do Ensino Fundamental e Médio, respeitadas as diversidades regional, estadual e local ( BRASIL, 2014 ).
Portanto, todos esses marcos regulatórios possibilitam à BNCC primar pelos princípios éticos, políticos e estéticos que visam à formação humana integral e à construção de uma sociedade justa, democrática e inclusiva, como fundamentado nas \textbf{Diretrizes} Curriculares Nacionais da Educação Básica ( DCN ).
De igual modo à BNCC, o \textbf{Currículo} do Estado de Alagoas está fundamentado no Art. 3º da LDB nº 9.394/96, nas \textbf{Diretrizes} Curriculares Nacionais, nos PCNs, na Lei nº 13.005/2014⁷, que promulgou o Plano Nacional de Educação ( PNE ), na Constituição de 1988, \textbf{Artigo} 205, e na Lei nº 7.795, de 22 de \textbf{janeiro} de 2016, que \textbf{instituiu} o Plano Estadual de Educação de Alagoas. 3.
A Área de Ciências Humanas Busca -se com a organização de um \textbf{caderno} para o componente da Área de Ciências Humanas a consolidação de uma proposta de trabalho coletivo e interdisciplinar.
Disto compreende -se que um \textbf{Currículo} para a área de Ciências Humanas pressupõe politicamente nos vislumbrarmos frente aos problemas sociais, econômicos, políticos, ambientais e culturais pelos quais passam o planeta, nosso país, nosso Estado e, respectivamente, cada \textbf{município}.
O \textbf{currículo} das Ciências Humanas, trabalhado no Ensino Fundamental Anos Iniciais e Finais, tem por finalidade contribuir para o desenvolvimento de sujeitos autônomos, conscientes e capazes de conviver e respeitar as diferenças.
Os componentes curriculares de Geografia e História, ao estimular os estudantes a desenvolverem uma compreensão do mundo, não favorecem apenas seu desenvolvimento autônomo, mas os tornam aptos a uma intervenção mais responsável no mundo em que vivem.
Nesse sentido, as Ciências Humanas devem, assim, estimular uma formação ética, elemento fundamental para a formação das novas gerações, auxiliando os estudantes a construir um sentido de responsabilidade para valorizar : os direitos humanos ; o respeito ao ambiente e à própria coletividade ; o \textbf{fortalecimento} de valores sociais, tais como a solidariedade, a participação e o protagonismo voltados para o bem comum ; e, sobretudo, a preocupação com as desigualdades sociais.
A BNCC orienta que : > Os conhecimentos específicos na área de Ciências Humanas exigem clareza na definição de um conjunto de objetos de conhecimento que favoreçam o desenvolvimento de habilidades e que aprimorem a capacidade de os estudantes pensarem diferentes culturas e sociedades, em seus tempos históricos, territórios e paisagens ( compreendendo melhor o Brasil, sua diversidade regional e territorial ).
E também que os levem a refletir sobre sua inserção singular e responsável na história da sua família, comunidade, nação e mundo ( 2017, p. 352 ).
Dessa forma, ela permite aos educandos conhecimentos que lhes possibilitem conviver em sociedade participando das tomadas de decisões de forma crítica nas questões sociais, éticas e políticas, condições basilares para uma atuação crítica e orientada por valores democráticos.
No Ensino Fundamental, os procedimentos de investigação em Ciências Humanas devem propiciar nos estudantes a capacidade de observação.
Na Geografia e na História, espera -se nessa etapa de ensino que seja trabalhado “ o reconhecimento do Eu e o sentimento de pertencimento dos estudantes à vida da família e da comunidade ” ( BNCC, 2017, p. 209 ).
As competências e habilidades propostas neste documento dizem respeito à formação para o exercício da cidadania.
Por isso, é essencial que os estudantes compreendam que são sujeitos capazes de interferir na realidade em que vivem e que os conhecimentos adquiridos na escola podem contribuir para seu desenvolvimento intelectual e social.
Cabe ainda às Ciências Humanas cultivar a formação de estudantes intelectualmente autônomos, com capacidade de articular categorias de pensamento histórico e geográfico em face de seu próprio tempo, percebendo as experiências humanas e refletindo sobre elas, com base na diversidade de pontos de vista.
A educação escolar deve propiciar, além da transmissão sistemática dos conteúdos de ensino, historicamente produzidos e acumulados, assegurar que os educandos se apropriem desses conteúdos de forma ativa, para que possam reelaborar esses conhecimentos e, com isso, obter um senso crítico mais concreto, embasado na compreensão científica e tecnológica da realidade social e política na qual estão inseridos socialmente.
Um dos grandes problemas da Geografia do século XXI é estimular o estudante a sua criatividade, sua imaginação, sua curiosidade, o raciocínio do pensamento geográfico e, consequentemente, levá -lo a compreender que o espaço geográfico ao longo dos tempos foi sendo modificado e transformado pela ação humana, e que, acima de tudo, seja capaz de tornar esse estudante um cidadão autônomo e crítico para atuar e mediar conflitos.
Cabe às Ciências Humanas cultivar a formação de estudantes intelectualmente autônomos, com capacidade de articular categorias de pensamento histórico e geográfico em face de matéria a ser ensinada e aprendida, isto impediria que ela dialogasse com outras disciplinas.
A História é uma ciência viva, construída pelos homens nos mais variados estádios da sociedade.
O componente de História apresentado aqui neste documento apresenta -se como uma ciência que analisa o tempo, em suas várias idades e períodos cronológicos.
Sendo assim, o seu objeto de estudo é a relação do presente com o passado no decorrer do tempo histórico das sociedades, buscando contribuir com a formação do cidadão autônomo e integral.
Portanto, os conhecimentos específicos abordados na Área de Ciências Humanas no Ensino Fundamental dialogam diretamente com as 10 Competências Gerais da BNCC e com as Competências Específicas de cada componente curricular Geografia/História.
Assim, eles devem contribuir para a ampliação de conhecimento sobre o mundo social e as questões políticas, éticas e sociais e desenvolver nos educandos a autonomia intelectual, princípio fundamental para atuação reflexiva e aquisição de valores democráticos.
O quadro nº 01 mostra as Competências Específicas da Área de Ciências Humanas para o Ensino Fundamental : 1.
Compreender a si e ao outro como identidades diferentes, de forma a exercitar o respeito à diferença em uma sociedade plural e promover os direitos humanos. 2.
Analisar o mundo social, cultural e digital e o meio técnico-científico-informacional com base nos conhecimentos das Ciências Humanas, considerando suas variações de significado no tempo e no espaço, para intervir em situações do cotidiano e se posicionar diante de problemas do mundo contemporâneo. 3.
Identificar, comparar e explicar a intervenção do ser humano na natureza e na sociedade, exercitando a curiosidade e propondo ideias e ações que contribuam para a transformação espacial, social e cultural, de modo a participar efetivamente das dinâmicas da vida social. 4.
Interpretar e expressar sentimentos, crenças e dúvidas com relação a si mesmo, aos outros e às diferentes culturas, com base nos instrumentos de investigação das Ciências Humanas, promovendo o acolhimento e a valorização da diversidade de indivíduos e de grupos sociais, seus saberes, identidades, culturas e potencialidades, sem preconceitos de qualquer natureza. 5.
Comparar eventos ocorridos simultaneamente no mesmo espaço e em espaços variados, e eventos ocorridos em tempos diferentes no mesmo espaço e em espaços variados. 6.
Construir argumentos, com base nos conhecimentos das Ciências Humanas, para negociar e defender ideias e opiniões que respeitem e promovam os direitos humanos e a consciência socioambiental, exercitando a responsabilidade e o protagonismo voltados para o bem comum e a construção de uma sociedade justa, democrática e inclusiva. 7.
Utilizar as linguagens cartográfica, gráfica e iconográfica e diferentes gêneros textuais e tecnologias digitais de informação e comunicação no desenvolvimento do raciocínio espaço-temporal relacionado à localização, distância, direção, duração, simultaneidade, sucessão, ritmo e conexão.

% matched lemmas: artigo, caderno, currículo, diretriz, fortalecimento, instituir, instituição, janeiro, município, parecer, política, prever, reiterar
\end{textsample}
